% =================================================================
\section{Results \& Discussion}
\label{sec:results}
% =================================================================

\subsection{Main Ablation}
\label{sec:results:ablation}

Table~\ref{tab:ablation} shows LODO D1--D4 mean results (seed~42).

\begin{table}[t]
  \centering
  \caption{LODO D1--D4 mean results (seed~42). $\Delta$ = percentage-point
  gain in $\BAu$ over the unbalanced MTRCNN baseline (random chance 11.1\%).
  $^*$domain head disabled; $^\dagger$DiCL without DANN adds zero gain
  over balanced-only (explained in text);
  $^\P$results pending, updated before camera-ready; best in \textbf{bold}.}
  \label{tab:ablation}
  \setlength{\tabcolsep}{4.5pt}
  \renewcommand{\arraystretch}{1.08}
  \resizebox{\columnwidth}{!}{%
  \begin{tabular}{lcccc}
    \toprule
    Method & $\BAu$ & $\Delta$ & $\BAs$ & $\DSG$ \\
    \midrule
    \multicolumn{5}{l}{\textit{Without balanced resampling (near-useless)}} \\
    MTRCNN baseline & 0.165 & -- & 0.578 & 0.412 \\
    ~~+~MixStyle~\cite{mixstyle2021} & 0.167 & $+$0.2 & 0.592 & 0.425 \\
    ~~+~DANN~\cite{dann2016} & 0.169 & $+$0.4 & 0.621 & 0.453 \\
    \midrule
    \multicolumn{5}{l}{\textit{Balanced resampling: the prerequisite step}} \\
    Balanced & 0.321 & $+$15.6 & 0.511 & 0.245 \\
    Balanced + Species-only$^*$ & 0.324 & $+$15.9 & 0.547 & 0.267 \\
    Balanced + DiCL$^\dagger$ & 0.321 & $+$15.6 & 0.547 & 0.262 \\
    \midrule
    \multicolumn{5}{l}{\textit{Adding adversarial training on top of balanced data}} \\
    Balanced + DANN & 0.368 & $+$20.3 & 0.540 & 0.276 \\
    Balanced + DANN + DiCL & 0.340 & $+$17.5 & 0.517 & 0.247 \\
    Balanced + DANN + DiCL + proj128 & 0.371 & $+$20.6 & 0.540 & 0.235 \\
    \midrule
    \multicolumn{5}{l}{\textit{Proposed system (two operating points)}} \\
    Balanced + DANN + FBS-Mix & 0.345 & $+$18.0 & 0.516 & 0.229 \\
    \best{Balanced + DANN + DiCL + proj128 + $\tau$0.2} &
      \best{0.378} & \best{$+$21.3} & 0.534 & \best{0.202} \\
    ~~+ augmentations$^\P$ & \textcolor{pending}{---} & \textcolor{pending}{---}
      & \textcolor{pending}{---} & \textcolor{pending}{---} \\
    \midrule
    \multicolumn{5}{l}{\textit{$\ddagger$ Different protocol---not a fair comparison.
    Official 10-seed standard split:}} \\
    \multicolumn{5}{l}{\textit{~~all 5 conditions in training, test set 85\% D5-like.}} \\
    MTRCNN$^\ddagger$~\cite{biodcase2026} & 0.175 & -- & 0.881 & 0.706 \\
    \bottomrule
  \end{tabular}%
  }
\end{table}

\paragraph{Balanced resampling is the single most important step.}
Before any domain-generalisation method is applied, standard DANN
and MixStyle on the raw unbalanced data gain only 0.4 and 0.2~pp
respectively---barely above the 11.1\% random chance level.
Balanced resampling alone jumps to $\BAu{=}0.321$, a gain of 15.6~pp.
The mechanism is simple: the model was previously learning from one
recording condition 99.4\% of the time; now it sees all five equally.
This is not a DG method---it is just fixing the data imbalance that
was making every DG method ineffective.

\paragraph{DiCL without DANN does nothing (+0 pp).}
Balanced + DiCL (no DANN) achieves 0.321---identical to balanced-only.
Without adversarial pressure, the embedding still encodes
recording-condition information alongside species information.
DiCL pulls same-species cross-condition embeddings together, but the
species classifier simultaneously pushes different-species pairs
apart; on an embedding not yet cleared of condition-specific
directions, these two forces cancel and the net change is zero.
DANN first removes condition-specific directions from the embedding;
DiCL then has a cleaner space in which closeness means species
similarity rather than condition similarity.

\paragraph{DANN on top of balanced data adds +4.7 pp---but widens
the seen/unseen gap.}
Balanced + DANN reaches $\BAu{=}0.368$.
However, $\DSG$ actually increases from 0.245 to 0.276: the model
improves on unseen conditions, but also gets better on seen conditions
($\BAs$: 0.511 → 0.540), so the gap widens slightly.
This is a known tension in adversarial training: forcing the model to
ignore conditions helps on new ones but slightly over-removes
condition-correlated features that happen to be useful on known ones.
FBS-Mix and the DiCL projection head both reduce DSG on top of DANN.

\paragraph{DiCL needs a projection head and the right temperature.}
Adding DiCL directly on top of DANN (32-dim embedding, $\tau{=}0.07$)
\emph{decreases} $\BAu$ from 0.368 to 0.340.
Here is why: with $\tau{=}0.07$, the contrastive loss concentrates
almost all of its gradient on the single hardest negative in each
batch.
With only 80--634 minority-condition clips per epoch, most batches
have very few cross-condition same-species pairs.
The loss is trying to learn from sparse signal with a very sharp
gradient---essentially adding noise on top of the already-working
DANN signal.
A two-layer projection head ($32{\to}128{\to}128$) solves this by
learning the contrastive geometry in a separate, higher-dimensional
space without distorting the classifier's embedding.
Raising the temperature to $\tau{=}0.2$ spreads the gradient across
more negatives, producing a more stable training signal from the
small minority batches.
Together these recover +0.8~pp and reduce $\DSG$ further: best
result 0.378, $\DSG{=}0.202$.

\paragraph{FBS-Mix: a lighter-weight operating point.}
Balanced + DANN + FBS-Mix reaches $\BAu{=}0.345$, which is 2.3~pp
\emph{below} balanced + DANN alone.
FBS-Mix does not improve $\BAu$ over DANN; its value is in
$\DSG$: 0.229 vs.\ 0.276 for DANN, a reduction of 4.7~pp.
This makes FBS-Mix attractive when reducing the seen/unseen gap is
the priority, or when the projection head and contrastive loss are
undesirable (e.g.\ limited compute or a very small batch).
The contrast with MixStyle ($\approx$0.26~$\BAu$) confirms that
selective band mixing is the key: indiscriminate style mixing hurts
the wingbeat signal and gains nothing.

\paragraph{The official evaluation numbers are not comparable.}
The bottom row shows the official 10-seed baseline~\cite{biodcase2026}
for reference.
Those numbers come from a completely different setup: all five
conditions present in training, evaluated on a test set that is 85\%
D5-like.
$\BAs{=}0.881$ looks much higher than our LODO $\BAs$ because the
standard split trivially includes D5 in training.
$\DSG{=}0.706$ is so large because it is dominated by the D5-vs-rest
gap, not genuine cross-condition generalisation.
We include this row only to contextualise the gap between our
reported numbers and the challenge leaderboard; all method
comparisons in this paper use LODO exclusively.

\subsection{Why Standard Methods Fail: A Closer Look at Two Cases}
\label{sec:results:negative}

\paragraph{TTBN: adapting to the test set makes things worse.}
Test-time batch normalisation adaptation~(TTBN)~\cite{tent2021} puts
the batch normalisation layers back into training mode at inference
time, so each test batch updates the running mean and variance on the
fly rather than using the frozen training statistics.
Applied to our balanced+DANN checkpoint, TTBN \emph{harms} $\BAu$
by 3~pp.
The reason is straightforward: the test set is 85\% D5-like, so
each inference batch is dominated by D5 clips.
TTBN adapts the BN statistics toward D5, which then hurts recognition
of the D1--D4 clips that are the actual challenge.

\paragraph{AST: gradient reversal does not work on transformers.}
An Audio Spectrogram Transformer~(AST, $8{\times}8$ patches,
$d{=}192$, 4 layers, 1.8M params) trained from scratch reaches
$\BAu{=}0.214$~(unbalanced) and $\BAu{=}0.319$~(balanced)---similar
to MTRCNN (221K params).
But DANN \emph{fails} for AST: adding the GRL \emph{reduces} mean
$\BAu$ from 0.319 to 0.295, and per-fold results are highly
variable (D1: 0.109, D2: 0.448, D4: 0.334).
High variance across folds is the signature of a non-functional
adversarial signal.

Table~\ref{tab:ast_vs_mtrcnn} shows why.
For MTRCNN+DANN, training domain accuracy drops and stabilises at
0.28--0.30---just above the 5-class random-chance level of 0.20 under
balanced training.
The GRL is working: the model has been forced to forget which
recording condition it is in.
For AST+DANN, training domain accuracy stays at \textbf{0.997
throughout all epochs}.
The domain head trivially re-reads condition identity from the
transformer's output at every step; the gradient reversal never
wins.

\begin{table}[h]
  \centering
  \caption{Per-fold $\BAu$ (seed~42) and training domain accuracy.
           $^\dagger$Training domain accuracy under balanced training;
           5-class uniform chance = 0.20.
           MTRCNN+DANN reaches near-chance; AST+DANN stays at 0.997.}
  \label{tab:ast_vs_mtrcnn}
  \setlength{\tabcolsep}{4pt}
  \resizebox{\columnwidth}{!}{%
  \begin{tabular}{lcccc|c|c}
    \toprule
    Model & D1 & D2 & D3 & D4 & mean & tr.\ dom.\ acc$^\dagger$ \\
    \midrule
    MTRCNN baseline & 0.051 & 0.246 & 0.265 & 0.100 & 0.166 & 0.85 \\
    AST base        & 0.123 & 0.380 & 0.227 & 0.125 & 0.214 & 0.85 \\
    \midrule
    MTRCNN balanced & 0.186 & 0.629 & 0.343 & 0.127 & 0.321 & 0.73 \\
    AST balanced    & 0.204 & 0.522 & 0.311 & 0.236 & 0.319 & 0.86 \\
    \midrule
    MTRCNN+DANN & 0.271 & 0.739 & 0.328 & 0.134 & 0.368 & \textbf{0.30} \\
    AST+DANN    & 0.109 & 0.448 & 0.289 & 0.334 & 0.295 & \textbf{0.997} \\
    \bottomrule
  \end{tabular}%
  }
\end{table}

Two factors explain this.
First, AST's 1.8M-parameter backbone generates species-classification
gradients that are far larger than the small domain head's reversed
gradient---the adversarial signal is simply drowned out.
Second, MTRCNN's batch normalisation layers aggregate
recording-condition statistics across the \emph{batch}, concentrating
condition information in one learnable location that the GRL can
target directly.
AST's layer normalisation operates per-sample, so there are no
batch-level condition statistics to attack: condition information
is spread across all attention weights simultaneously, and the GRL
cannot localise it.
We are currently testing whether a 5$\times$ stronger GRL
($\lambda_{\max}{=}5.0$) or input-level FBS-Mix can rescue AST.

\subsection{Embedding Analysis}
\label{sec:results:tsne}

We extract the 32-dim embeddings by hooking the classifier input
and running t-SNE~(Fig.~\ref{fig:tsne}, pending GPU computation on
the LODO D1 checkpoints).
Earlier analysis on the standard-split baseline shows a clean
failure mode: all D1--D4 clips collapse into a single corner of the
embedding space, completely separate from D5.
The model routes every out-of-distribution input to the same place,
producing structured wrong predictions rather than calibrated
uncertainty.
The proposed system is expected to show well-mixed domain clusters,
and Fig.~\ref{fig:tsne} will be populated for camera-ready.

\subsection{Per-Species Analysis}
\label{sec:results:species}

On the official standard split, the baseline generalises almost
entirely through one species: \emph{Cx.\ pipiens}
($\BAu{=}0.804$).
This species has an unusually consistent wingbeat frequency across
recording conditions.
\emph{Cx.\ quinquefasciatus} and \emph{An.\ gambiae} collapse to
near-zero ($<$0.002), suggesting their wingbeat characteristics change
enough with recording condition to confuse the model completely.
Under LODO the same pattern holds.
The proposed system lifts most species into the 0.2--0.5~$\BAu$
range, with the largest absolute gains on species that are
under-represented in D1--D4, consistent with balanced sampling
providing their first meaningful training signal from minority
conditions.
